\documentclass[leqno]{article}

%% Created with wxMaxima 17.05.0

\setlength{\parskip}{\medskipamount}
\setlength{\parindent}{0pt}
\usepackage[utf8]{luainputenc}
\DeclareUnicodeCharacter{00B5}{\ensuremath{\mu}}
\usepackage{graphicx}
\usepackage{color}
\usepackage{amsmath}
\usepackage{ifthen}
\usepackage{breqn}   % don't for get this line
\newsavebox{\picturebox}
\newlength{\pictureboxwidth}
\newlength{\pictureboxheight}
\newcommand{\includeimage}[1]{
    \savebox{\picturebox}{\includegraphics{#1}}
    \settoheight{\pictureboxheight}{\usebox{\picturebox}}
    \settowidth{\pictureboxwidth}{\usebox{\picturebox}}
    \ifthenelse{\lengthtest{\pictureboxwidth > .95\linewidth}}
    {
        \includegraphics[width=.95\linewidth,height=.80\textheight,keepaspectratio]{#1}
    }
    {
        \ifthenelse{\lengthtest{\pictureboxheight>.80\textheight}}
        {
            \includegraphics[width=.95\linewidth,height=.80\textheight,keepaspectratio]{#1}
            
        }
        {
            \includegraphics{#1}
        }
    }
}
\newlength{\thislabelwidth}
\DeclareMathOperator{\abs}{abs}
\usepackage{animate} % This package is required because the wxMaxima configuration option
                      % "Export animations to TeX" was enabled when this file was generated.

\definecolor{labelcolor}{RGB}{100,0,0}
\begin{document}
\begin{minipage}[t]{8ex}\color{red}\bf
(\% i13)
\end{minipage}
\begin{minipage}[t]{\textwidth}\color{blue}
p7:subst(pt, P, p4);
\end{minipage}
%%%%%%%%%%%%%%
%%% OUTPUT %%% 'show full equation,*.pdf'
%%%%%%%%%%%%%%
%%%%%%%%%%%%%%%%%%%%%%%%%%%%%%%%%%%%%%%%%%%%%%%%
% 1e) replace :  
% =>   \[\displaystyle -> \begin{dmath}                  see above '\usepackage{breqn}
% 2e) replace :  
% =>  \tag{p7}         ->  p7]
% 3e) remove :  
% =>  \mbox{}
% 4e) replace :
% =>  \]               -> \end{dmath}
%%%%%%%%%%%%%%%%%%%%%%%%%%%%%%%%%%%%%%%%%%%%%%%%
\begin{dmath}
p7]
\mathit{Z2}=\mathit{B1}\, \left( \frac{R T}{V}+\frac{B R T}{{{V}^{2}}}+\frac{C R T}{{{V}^{3}}}+\frac{D R T}{{{V}^{4}}}\right) +\mathit{D1}\,
 {{\left( \frac{R T}{V}+\frac{B R T}{{{V}^{2}}}+\frac{C R T}{{{V}^{3}}}+\frac{D R T}{{{V}^{4}}}\right) }^{3}}+\mathit{C1}\, 
{{\left( \frac{R T}{V}+\frac{B R T}{{{V}^{2}}}+\frac{C R T}{{{V}^{3}}}+\frac{D R T}{{{V}^{4}}}\right) }^{2}}+1
\end{dmath}
\end{document}
